\documentclass[12pt,a4paper]{ctexart}
\usepackage{amsmath,amssymb}
\usepackage{booktabs}
\usepackage{array}
\usepackage{geometry}
\usepackage{enumitem}
\usepackage[hidelinks]{hyperref}
\geometry{margin=2.5cm}
\setlength{\parskip}{0.5em}
\setlength{\parindent}{2em}
\linespread{1.25}

\title{\bfseries 2026 CUMCM A 题 参考文献清单\\（教材 + 论文，均经摘要核实）}
\author{检索时间：2026 年 9 月}
\date{}

\begin{document}
\maketitle

\begin{abstract}
来源：Web 检索（Google/Bing 索引）、Semantic Scholar、ScienceDirect/RSC/MDPI/CNKI 条目页的摘要级内容。标记说明：\textbf{[摘要已核]} $=$ 已读到该文摘要原文/摘要级内容并确认与本题相关；\textbf{用途} $=$ 建议用在论文哪一部分。
注：本机网络策略禁止直连部分出版社网站（pubs.rsc.org、sciencedirect 均被拦截），以下摘要核实基于搜索引擎返回的摘要文本；引用前建议用学校图书馆数据库（知网/Web of Science/Elsevier/RSC）再核对一遍原文。

\textbf{符号说明}（与《A题问题1知识点串讲》符号表一致）：$D$——水分扩散系数，m$^2$/s；$D_{\mathrm{eff}}$——有效水分扩散系数（文献中由干燥曲线回归得到的等效值），m$^2$/s；$E_a$——表观活化能，kJ/mol；$M$——水分含量（各文献定义略有差异，按原文语境理解）；$T$——温度，$^\circ$C；$C$——水分浓度（干基含水率），kg/kg；$r$——径向坐标，m；$t$——时间，s；$h$——对流换热系数，W/(m$^2\cdot$K)；$h_m$——对流传质系数，m/s。
\end{abstract}

\tableofcontents
\newpage

\section{教材与专著（打基础用，按阅读顺序）}

\subsection{Crank, J. \textit{The Mathematics of Diffusion}, 2nd ed., Clarendon Press / Oxford University Press, 1975（ISBN 9780198534112）\quad 本题圣经}
\begin{itemize}
\item 覆盖：扩散方程、圆柱中的扩散（第 5 章，含 Bessel 级数解）、\textbf{移动边界}（第 7 章）、有限差分法（第 8 章）、变扩散系数、\textbf{热与水分同时扩散}（末章）。
\item 用途：问题 1--4 的方程、解析验证锚点（冻结 $D$ 时的级数解）、移动边界处理，全部能在本书找到对应章节。有中译本《扩散的数学》。
\item 姊妹篇：Carslaw, H.S. \& Jaeger, J.C. \textit{Conduction of Heat in Solids}, 2nd ed., Oxford, 1959——热传导侧的同款工具书，两书互相引用，合起来是本题的标准理论底座。
\end{itemize}

\subsection{Incropera, F.P., DeWitt, D.P., Bergman, T.L., Lavine, A.S. \textit{Fundamentals of Heat and Mass Transfer}, 6th ed. (Wiley, 2006) / 8th ed. (Wiley, 2017)}
\begin{itemize}
\item 第 5 章 Transient Conduction：5.4--5.6 \textbf{圆柱径向瞬态导热与对流边界}（5.6.1 精确解、5.6.2 近似解）、5.10 有限差分（显式/隐式）；附录 D Heisler 图。
\item 第 14 章 Diffusion Mass Transfer：Fick 定律、14.7 \textbf{瞬态扩散}（与第 5 章完全平行的类比）。
\item 用途：模块 A/B 的标准教材，第 5 章 5.6 节可直接对着写问题 1 的线性部分。
\end{itemize}

\subsection{Whitaker, S. ``Simultaneous Heat, Mass, and Momentum Transfer in Porous Media: A Theory of Drying'', \textit{Advances in Heat Transfer}, Vol. 13, pp. 119--203, 1977（DOI 10.1016/S0065-2717(08)70223-5）\quad 干燥理论奠基文献}
\begin{itemize}
\item 体积平均法建立多孔介质干燥的输运方程组；已被引 1183+ 次。1998 年 Whitaker 有续篇 ``Coupled Transport in Multiphase Systems: A Theory of Drying''（\textit{Advances in Heat Transfer}, pp. 1--102）。
\item 用途：论文「问题重述/模型假设」部分引用，说明为什么可以用连续介质扩散-导热框架描述药材（多孔介质）；也是「为什么题给经验式就够用」的理论台阶。
\end{itemize}

\subsection{Mujumdar, A.S. (ed.) \textit{Handbook of Industrial Drying}, 4th ed., CRC Press, 2015（ISBN 9781466596658）}
\begin{itemize}
\item 相关章节：Ch.1 Principles/Classification/Selection of Dryers（干燥速率曲线、含水率概念）、``Basic Process Calculations and Simulations in Drying''、``Transport Properties in the Drying of Solids''。
\item 用途：干燥三阶段、干基/湿基、临界/平衡含水率的权威出处；引言部分引用。
\end{itemize}

\subsection{杨世铭, 陶文铨. 《传热学》（第 4/5 版）, 高等教育出版社}
\begin{itemize}
\item 第 2--3 章稳态/非稳态导热（集总参数法、毕渥数、无限长圆柱）、第 4 章数值解法。
\item 用途：中文快速入门传热模块 A（比 Incropera 薄、中文好读）。
\end{itemize}

\subsection{化工原理教材（任选一本，看「干燥」一章）：陈敏恒《化工原理》下册 / 柴诚敬《化工原理》}
\begin{itemize}
\item 干燥章：湿空气性质、干燥速率曲线（预热/恒速/降速）、临界与平衡含水率、物料衡算。
\item 用途：模块 C 全部概念；中文教材例子最贴合本题背景。
\end{itemize}

\section{直接相关论文（几何与模型和本题几乎一致）——精读}

\subsection{P1. Simal, S., Rossell\'o, C., Berna, A., Mulet, A. ``Drying of shrinking cylinder-shaped bodies'', \textit{Journal of Food Engineering}, 37(4): 423--435, 1998（DOI 10.1016/S0260-8774(98)00095-8）\quad [摘要已核]}
\begin{itemize}
\item 摘要要点：圆柱固体干燥过程中温度、平均水分与水分分布的预测模型；\textbf{考虑收缩}；宏观热平衡$+$菲克定律微观质平衡同时求解（Runge-Kutta-Merson $+$ 有限差分）；\textbf{有效扩散系数写成温度与局部含水率的函数}；用 90\,$^\circ$C 西兰花茎干燥曲线标定，解释方差 99.8\%。
\item 用途：问题 4（收缩）的\textbf{模板论文}，其「热平衡$+$扩散$+$移动边界$+D(T,C)$」结构就是本题问题 2--4 的骨架；模型假设写法（如「导热系数视为无限大$\Rightarrow$温度均匀」）值得对比讨论。
\end{itemize}

\subsection{P2. Abbasi Souraki, B., Mowla, D. ``Axial and radial moisture diffusivity in cylindrical fresh green beans in a fluidized bed dryer with energy carrier: Modeling with and without shrinkage'', \textit{Journal of Food Engineering}, 88(1): 9--19, 2008（DOI 10.1016/j.jfoodeng.2007.05.013）\quad [摘要已核]}
\begin{itemize}
\item 摘要要点：圆柱绿豆用\textbf{无限/有限圆柱菲克模型}求轴向与径向有效扩散系数；\textbf{含收缩与不含收缩}两种建模对比；收缩时扩散系数估计值更小；轴向 $D$ 显著大于径向 $D$；$D$ 与温度呈 \textbf{Arrhenius} 关系。
\item 用途：问题 1 圆柱几何「1D 径向模型」合理性的佐证与反例讨论（长径比大$\Rightarrow$径向主导）；「忽略收缩会高估 $D$」可直接引用到问题 4 动机。
\end{itemize}

\subsection{P3. ``Heat and moisture transport in okra cylinders with shrinkage effects under solar drying: a multiphysics-based simulation approach'', \textit{Sustainable Food Technology} (RSC), 2025（DOI 10.1039/D4FB00343H）\quad [摘要已核]}
\begin{itemize}
\item 要点：秋葵\textbf{圆柱}（半径 6\,mm）干燥的\textbf{热-湿耦合}瞬态模拟；2D 轴对称；COMSOL 耦合「固体传热」$+$「稀物质传递」；\textbf{ALE 动网格处理收缩}；表面对流干燥边界（含蒸发项）。
\item 用途：最现代化的同类算例（2025 年），问题 4 的 ALE/移动网格实现细节、边界条件写法都能对照；也可作为「多物理场软件 vs 自编 FDM」的方法对比引用。
\end{itemize}

\subsection{P4. 任迪峰.《中药材干燥过程中质量退化及优化干燥工艺的研究》（博士学位论文，中国农业大学）}
\begin{itemize}
\item 要点：第四章把常见根类中药材假设为\textbf{近似变径轴对称圆柱体}，建立温度传导与水分扩散\textbf{耦合的二维传热传质有限元模型}（轴对称圆柱坐标微元体热量/质量平衡）。
\item 用途：\textbf{中文文献里与本题重合度最高的一份}——同样是「根类药材$+$轴对称圆柱$+$热质耦合 FEM」；其文献综述还梳理了 Chau、Irudayaraj、Liu、Sakai、Hayakawa 等一维/二维热质耦合建模的谱系，可直接当「文献综述」底稿。
\end{itemize}

\subsection{P5. Kim, M.H., Kim, C.S., Park, S.J., Lee, C.H. ``Numerical Analysis of Moisture Diffusion for Hot-air Drying of Ginseng''（人参热风干燥水分扩散的数值解析研究），《韩国农业机械学会志》, 1998}
\begin{itemize}
\item 要点：人参视为\textbf{圆柱体}，水分仅径向迁移，控制方程
$\dfrac{\partial M}{\partial t}=\dfrac{1}{r}\dfrac{\partial}{\partial r}\Big[D_{\mathrm{eff}}(T,M)\,r\,\dfrac{\partial M}{\partial r}\Big]$，
数值求解抛物线 PDE 模拟热风干燥。
\item 用途：问题 1 水分方程（含 $D(C)$）的\textbf{一字不差的先例}；1999 年同团队在 \textit{Food Engineering Progress} 还有缓苏（tempering）干燥的隐式有限差分续作，可作为时间步长/隐式格式选择参照。
\end{itemize}

\subsection{P6. 庞凌云, 袁志华, 詹丽娟, 等. 《铁棍山药片热风干燥过程中传热传质规律》，《食品与机械》40(8), 2024（DOI 10.13652/j.spjx.1003.5788.2023.60171）}
\begin{itemize}
\item 要点：山药片按\textbf{圆柱体}建模（半径约 20\,mm），测定导热系数/比热容，ANSYS Transient Thermal 模拟温度场；Modified Page 模型拟合传质；中心温度最低、外表面最高，内外温差由 7.43\,K 降至 5.82\,K。
\item 用途：中文写作范例（结构、术语）；「圆柱药材热风干燥中心-表面温差」的实测佐证；问题 1 表 1 温度形态（外高内低、梯度渐小）的实验对应。
\end{itemize}

\subsection{P7. 王学成, 康超超, 伍振峰, 等. 《二至丸热风干燥过程温度均匀性模拟与实验》，《中草药》51(5): 1226--1232, 2020}
\begin{itemize}
\item 要点：COMSOL 传热传质模型$+$Fick 第二定律平板模型求 $D_{\mathrm{eff}}$（0.76--3.94$\times10^{-7}$\,m$^2$/s），60/80/100\,$^\circ$C 热风，探针测温验证。
\item 用途：中药材$+$COMSOL$+$Fick 的中文范式；$D_{\mathrm{eff}}$ 量级可与题给公式互证。
\end{itemize}

\subsection{P8. Pacheco-Aguirre, F.M., Ladr\'on-Gonz\'alez, A., Ruiz-Espinosa, H., Garc\'ia-Alvarado, M.A., Ruiz-L\'opez, I.I. ``A method to estimate anisotropic diffusion coefficients for cylindrical solids: Application to the drying of carrot'', \textit{Journal of Food Engineering}, 125: 24--33, 2014（DOI 10.1016/j.jfoodeng.2013.10.015）\quad [摘要已核]}
\begin{itemize}
\item 摘要要点：基于\textbf{有限各向异性圆柱}非稳态传质方程\textbf{解析解}（含 Bessel 函数）估计轴向/径向/角向扩散系数；胡萝卜 2.2\,cm 直径圆柱，80\,$^\circ$C；$D=0.53$--$2.93\times10^{-9}$\,m$^2$/s；三方向差异显著。
\item 用途：圆柱扩散解析解的现成公式库（验证你的 FDM 代码）；「各向异性是否必要」的讨论素材（本题长 25\,cm、径 2\,cm，长径比 12.5，可用它论证径向 1D 近似的合理性）。
\end{itemize}

\section{方法与理论支撑论文}

\subsection{P9. Chandra Mohan, V.P., Talukdar, P. ``Three dimensional numerical modeling of simultaneous heat and moisture transfer in a moist object subjected to convective drying'', \textit{International Journal of Heat and Mass Transfer}, 53(21--22): 4638--4650, 2010（DOI 10.1016/j.ijheatmasstransfer.2010.06.029）\quad [摘要已核]}
\begin{itemize}
\item 要点：耦合热湿方程 3D 数值解（有限体积法、全隐式）；$h$ 由自编 CFD 给出，$h_m$ 由\textbf{热-浓度边界层类比}（Chilton--Colburn）换算；气流 313$\to$353\,K 省时 40\%。
\item 用途：数值方法（隐式、耦合求解）写法；「变表面换热系数 vs 常数系数」的讨论——支持本题「常数 $h/h_m$」简化假设的合理性论证。
\end{itemize}

\subsection{P10. Defraeye, T., Blocken, B., Carmeliet, J. ``Analysis of convective heat and mass transfer coefficients for convective drying of a porous flat plate by conjugate modelling'', \textit{International Journal of Heat and Mass Transfer}, 55(1--3): 112--124, 2012（DOI 10.1016/j.ijheatmasstransfer.2011.08.047）\quad [摘要已核]}
\begin{itemize}
\item 要点：共轭建模（CFD 空气域$+$多孔材料域）反推对流系数，发现系数在恒速段/降速段近似常数、\textbf{在段间过渡处才显著变化}；热质类比在两段内成立、过渡处偏差大。
\item 用途：模型假设处引用「常数 $h$、$h_m$ 在常规对流干燥中可接受」，堵住评阅人「为什么不考虑系数时空变化」的疑问。
\end{itemize}

\subsection{P11. Garc\'ia del Valle, J., Sierra Pallares, J. ``Analytical solution for the coupled heat and mass transfer formulation of one-dimensional drying kinetics'', \textit{Journal of Food Engineering}, 230: 99--113, 2018（DOI 10.1016/j.jfoodeng.2018.02.029）\quad [摘要已核]}
\begin{itemize}
\item 要点：耦合热质守恒的\textbf{解析解}（内部解耦、边界耦合：内部毛细扩散$+$界面蒸发）；自然涌现\textbf{两个干燥阶段}（对流控制段$\to$扩散控制段）；假设恒定 $D_{\mathrm{eff}}$、无收缩。
\item 用途：理解本题「为什么边界条件是耦合的、内部方程是解耦/弱耦合的」；解析解可作数值代码验证锚点；论文「模型合理性」部分的引用。
\end{itemize}

\subsection{P12. Garc\'ia-Alvarado, M.A., Pacheco-Aguirre, F.M., Ruiz-L\'opez, I.I. ``Analytical solution of simultaneous heat and mass transfer equations during food drying'', \textit{Journal of Food Engineering}, 142: 39--45, 2014}
\begin{itemize}
\item 要点：Laplace 变换求食品干燥热质同时传递解析解，考虑温度对界面水分的影响。
\item 用途：与 P11 同族，问题 2 两阶段解析讨论备用。
\end{itemize}

\section{中药材干燥数据与综述（参数标定、量级互证、引言背景）}

\subsection{P13. 《中草药干燥加工现状及发展趋势》，《南京中医药大学学报》, 2021（DOI 10.14148/j.issn.1672-0482.2021.0786）——中文综述}
用途：引言「中药材干燥背景与难点」整段可直接取材；经验/半理论/理论模型分类（Weibull、Page、Midilli 等）适合在论文里交代「为什么不直接用经验模型、而要建理论模型」。

\subsection{P14. ``Drying kinetics of Salviae Miltiorrhizae Radix et Rhizoma (丹参) and dynamics of active components in drying process''（PubMed PMID 39929654, 2025）}
\begin{itemize}
\item 要点：丹参热风干燥 $D_{\mathrm{eff}}=1.746\times10^{-8}\sim9.354\times10^{-8}$\,m$^2$/s，活化能 \textbf{40.0\,kJ/mol}（红外 26.62）。
\item 用途：题给公式 $D=2.4\times10^{-3}e^{-0.45/C}e^{-3850/T}$ 的 $E_a\approx32$\,kJ/mol、$D\sim10^{-9}$--$10^{-8}$\,m$^2$/s 与实测完全同量级 $\Rightarrow$ \textbf{论文里用这类数据给题给经验式做「物理合理性背书」}（高分讨论点）。
\end{itemize}

\subsection{P15. 《鲜地黄干燥动力学及其对环烯醚萜类成分影响研究》（2024）}
要点：热风 40--60\,$^\circ$C，Midilli 模型最优；$D_{\mathrm{eff}}=6.681\times10^{-10}\sim8.969\times10^{-9}$\,m$^2$/s，$E_a=34.94$\,kJ/mol。用途：量级互证；「变温干燥（60\,$^\circ$C 烘 42\,h 后 40\,$^\circ$C）」可作问题 2/3 分段策略的工艺背景。

\subsection{P16. ``Kinetics and variation of volatile components of Atractylodis Macrocephalae Rhizoma (白术) during hot-air drying''，《中国中药杂志》}
要点：30--70\,$^\circ$C 热风，Midilli 模型；$D_{\mathrm{eff}}=1.04\times10^{-9}\sim6.28\times10^{-9}$\,m$^2$/s，$E_a=37.47$\,kJ/mol。用途：量级互证$+$根茎类药材（与本题圆柱形药材同型）干燥行为参考。

\subsection{P17. 李坤.《根茎类中药材热风干燥的失水动力学模型及特性机理研究》（硕士论文，云南师范大学, 2020）}
要点：根茎类药材（占植物类中药材 80\% 以上）热风干燥传热传质理论；基于质量/热量守恒建立失水动力学模型；计算 $D_{\mathrm{eff}}$ 与活化能；分析温度、风速、相对湿度、厚度的影响。用途：中文「理论模型」写作范例$+$根茎类药材参数影响规律。

\section{检索渠道与检索式（后续自己查补）}

\begin{itemize}
\item \textbf{Google Scholar}（scholar.google.com）检索式示例：
  \begin{itemize}
  \item \texttt{"drying" AND "cylinder" AND ("shrinkage" OR "moving boundary")}
  \item \texttt{"coupled heat and mass transfer" drying numerical simulation}
  \item \texttt{"effective moisture diffusivity" cylinder Arrhenius}
  \item 中文：\texttt{中药材 热风干燥 动力学 有效水分扩散系数}；\texttt{圆柱 干燥 数值模拟 传热传质}
  \end{itemize}
\item \textbf{Semantic Scholar}（semanticscholar.org）：按 P1、P2 的「引用/被引」顺藤摸瓜（P1 被引 83+，P9 被引 130+，Whitaker 1977 被引 1183+）。
\item \textbf{知网 CNKI}：中文论文与学位论文（P4、P6、P7、P13--P17）；检索式 \texttt{中药材 AND 干燥 AND (数值模拟 OR 传热传质 OR 动力学)}。
\item \textbf{Web of Science / Scopus}：英文权威回溯，按 DOI 定位。
\item 获取全文：学校图书馆数据库（Elsevier/RSC/Wiley/知网）均可下载上述文献；部分老文献（P1/P2/P5）也在 ResearchGate 有作者上传版。
\end{itemize}

\section{文献 $\leftrightarrow$ 论文章节映射（写作速查）}

\begin{center}
\begin{tabular}{p{4.6cm}p{10.6cm}}
\toprule
论文部分 & 建议引用 \\
\midrule
问题重述/背景 & P13 综述、Mujumdar 手册、P4 学位论文综述段 \\
模型假设（连续介质、常数 $h/h_m$、径向 1D） & Whitaker 1977、P10 Defraeye、P8（长径比论证）、P2（轴向 vs 径向） \\
控制方程与边界条件 & Crank、Incropera、P5（人参圆柱方程）、P1（热平衡$+$Fick 结构） \\
参数来源与合理性 & 附录经验式$+$P14--P16 实测 $D_{\mathrm{eff}}/E_a$ 量级互证 \\
数值方法 & Incropera 5.10、Crank 第 8 章、P9（隐式耦合求解） \\
模型验证 & P11/P12 解析解、Crank 圆柱级数解、P6（中心-表面温差形态） \\
问题 4 收缩 & P1（模板）、P3（ALE）、P2（收缩对 $D$ 的影响） \\
灵敏度/讨论 & P9（变系数）、P10（过渡段）、P6/P7（温度均匀性） \\
\bottomrule
\end{tabular}
\end{center}

\end{document}
